\section*{Chapitre \ref{ch:m1:futures-contracts-and-exchanges} --- Contrats à terme et leurs bourses}

\begin{solution}{exo:m1:futures-contracts-and-exchanges:1}
$(6\,012.50 - 5\,987.25)/0.25 = 101$ \omterm{def:m1:futures-contracts-and-exchanges:tick}{ticks}~; $101 \times \$12.50 \times 12 =
\$15\,150$ de perte.
\end{solution}

\begin{solution}{exo:m1:futures-contracts-and-exchanges:2}
$111 + 24.5/32 = 111.765625$ et $112 + 8/32 = 112.25$~: une hausse de $0.484375$
point, soit 31 \omterm{def:m1:futures-contracts-and-exchanges:tick}{ticks} de $1/64$, valant $31 \times \$15.625 = \$484.375$.
\end{solution}

\begin{solution}{exo:m1:futures-contracts-and-exchanges:3}
\omterm{def:m1:futures-contracts-and-exchanges:future}{Contrat à terme} EURO STOXX 50, échéance de mars (H) de la première année finissant par 7 qui n'est pas
passée~: mars 2027. Dernier jour de négociation~: le troisième vendredi, 19 mars 2027, à 12:00
CET. \texttt{ESH6} ne peut pas correspondre à mars 2026, qui est expiré~: lu aujourd'hui, il
désigne mars 2036, qui n'est pas coté. Un symbole dont l'année ne comporte qu'un chiffre est
ambigu sans date, raison pour laquelle les données de référence portent
les dates d'échéance complètes.
\end{solution}

\begin{solution}{exo:m1:futures-contracts-and-exchanges:4}
\omterm{def:m1:what-a-trading-firm-does:market-maker}{Teneur de marché}~: la fourchette minimale, un \omterm{def:m1:futures-contracts-and-exchanges:tick}{tick}, vaut 1.85 points de base contre
0.42, chaque aller-retour capturé rapporte donc davantage, et la file d'attente à chaque cours est
longue~: la priorité dans la file, donc la vitesse et l'algorithme d'appariement,
décident qui en profite. Les cours bougent moins souvent, si bien que les cotations restent sûres plus longtemps. Fonds~:
50~millions d'euros correspondent à 926 contrats~; traverser une fourchette d'un \omterm{def:m1:futures-contracts-and-exchanges:tick}{tick} coûte un
\omterm{def:m1:futures-contracts-and-exchanges:tick}{demi-tick}, $926 \times \text{\euro}5 = \text{\euro}4\,630$, soit 0.93 point de base, là où le
même échange dans un contrat à \omterm{def:m1:futures-contracts-and-exchanges:tick}{tick} plus fin coûterait environ un quart de
ce coût.
\end{solution}

\begin{solution}{exo:m1:futures-contracts-and-exchanges:5}
$300\,000/21\,000 = 14.3$ fois~; une variation de 7~\% fait perdre le dépôt. Le dépôt
est le minimum exigé aujourd'hui par la chambre de compensation. Les pertes doivent être
payées en espèces chaque jour au titre de la \omterm{def:m1:clearing-and-settlement:margin}{marge de variation}, et la \omterm{def:m1:clearing-and-settlement:margin}{marge initiale} elle-même est relevée quand
la volatilité augmente~: l'argent nécessaire pour \emph{tenir} la position pendant une mauvaise
semaine représente plusieurs fois le dépôt (\cref{ch:m1:margin}).
\end{solution}

\begin{solution}{exo:m1:futures-contracts-and-exchanges:6}
Après que A a acheté 5 à B~: A $+5$, B $-5$, \omterm{def:m1:futures-contracts-and-exchanges:oi}{positions ouvertes} 5. C achète 3 à A~: A
$+2$, C $+3$, B $-5$~; A a cédé 3 contrats à C, \omterm{def:m1:futures-contracts-and-exchanges:oi}{positions ouvertes} 5. B achète
2 à C~: B $-3$, C $+1$~; les deux positions sont clôturées, \omterm{def:m1:futures-contracts-and-exchanges:oi}{positions ouvertes} 3. D achète 4 à A~: A
$-2$, D $+4$~; A clôture 2 et ouvre une position vendeuse de 2, D ouvre 4~: \omterm{def:m1:futures-contracts-and-exchanges:oi}{positions ouvertes} 5
(acheteurs~: C 1, D 4~; vendeurs~: A 2, B 3). Volume 14, \omterm{def:m1:futures-contracts-and-exchanges:oi}{positions ouvertes} 5.
\end{solution}

\begin{solution}{exo:m1:futures-contracts-and-exchanges:7}
Résidu le plus élevé~: 14.29~\% avec le contrat de 300\,000~\$ (à 1.05~million de dollars,
un demi-contrat non couvert) et 1.00~\% avec le petit contrat. Le résidu est au
plus un demi-contrat~: $150\,000/V$ et $15\,000/V$, soit $14.3\%$ et
$1.4\%$ pour $V = 1.05$ million de dollars~; les données respectent les deux bornes.
\end{solution}

\begin{solution}{exo:m1:futures-contracts-and-exchanges:8}
Chaque contrat ouvert a un acheteur et un vendeur~: les acheteurs ne dépassent jamais
en nombre les vendeurs. Une hausse des \omterm{def:m1:futures-contracts-and-exchanges:oi}{positions ouvertes} signifie que de nouveaux acheteurs
\emph{et} de nouveaux vendeurs ont ouvert des positions~; le cours a monté parce que les
acheteurs étaient le côté le plus impatient, ce que le cours dit déjà. Ce que les données peuvent révéler~:
si un mouvement s'accompagne de nouvelles positions (hausse des \omterm{def:m1:futures-contracts-and-exchanges:oi}{positions ouvertes}) ou
d'une liquidation (baisse des \omterm{def:m1:futures-contracts-and-exchanges:oi}{positions ouvertes})~; où en est le roll~; et, combiné
au rapport hebdomadaire du régulateur, quelle catégorie de trader détient quel côté.
\end{solution}

\begin{solution}{pb:m1:futures-contracts-and-exchanges:1}
\textbf{1.} 300\,000~\$ et 54\,000~\euro.
\textbf{2.} $1.15 \times 187\,000\,000/300\,000 = 716.83$~: vendre 717. Résidu
$215.05 - 215.10 = -\$50\,000$~: surcouverture d'un sixième de contrat.
\textbf{3.} $0.90 \times 43\,000\,000/54\,000 = 716.67$~: vendre 717~; résidu
$-\text{\euro}18\,000$.
\textbf{4.} Décembre 2026, \texttt{ESZ6} et \texttt{FESXZ6}, dernier jour de négociation
vendredi 18 décembre 2026, dans exactement trois mois. Le E-mini se règle
sur la \emph{cotation d'ouverture} de ce jour, et le FESX à midi~: la
couverture prend fin quelques heures avant la clôture du dernier jour.
\textbf{5.} $717 \times 300\,000 \times 5\% = \$10\,755\,000$.
\textbf{6.} \omterm{def:m1:futures-contracts-and-exchanges:future}{Contrats à terme}~: $717 \times 300\,000 \times 2\% = +\$4\,302\,000$. Poche~:
$1.15 \times 2\% \times 187$ millions $= -\$4\,301\,000$.
\textbf{7.} Le gain sur les \omterm{def:m1:futures-contracts-and-exchanges:future}{contrats à terme} est versé en espèces le lendemain matin~; la perte sur
les actions est latente. Lors d'une bonne journée, c'est l'inverse~: le fonds doit
\emph{verser} 4.3~millions de dollars de \omterm{def:m1:clearing-and-settlement:margin}{marge de variation} en espèces contre un gain latent,
et doit disposer de ces liquidités ou les emprunter.
\textbf{8.} $4\,302\,000 - 2.9\% \times 187\,000\,000 = -\$1\,121\,000$~: la
poche a reculé de 0.6~\% de plus que ce que prévoyait son bêta. Il s'agit du risque idiosyncratique et
factoriel du portefeuille, que ne supprime pas une couverture par indice.
\textbf{9.} $0.10 \times 187$ millions $= 18.7$ millions de dollars d'exposition à l'indice,
374 fois le résidu d'arrondi.
\textbf{10.} $717 \times 0.05 \times \$50 = \$1\,792.50$.
\textbf{11.} La devise~: la poche est un actif en euros pour un fonds en dollars, et
les gains et pertes du FESX sont aussi en euros. La couverture supprime le
bêta actions, et non les 43~millions d'euros d'exposition de change (elle ajoute en fait une petite
exposition en euros sur ses marges et son P\&L).
\textbf{12.} Un petit bruit de moyenne nulle~: les deux cours sont relevés au
même instant par des mécanismes différents, une moyenne sur trente secondes pour le \omterm{def:m1:futures-contracts-and-exchanges:future}{contrat
à terme} et un seul cours d'enchères pour chaque action. Les jours ordinaires, il s'agit d'un point de base
ou deux~; les jours de rééquilibrage d'indice et d'échéance, quand l'enchère de clôture
bouge, l'écart est plus grand, et il s'inverse le lendemain.
\textbf{13.} Liquidité et coût~: le E-mini est la couverture la moins chère par dollar.
Il renonce à la correspondance~: une poche à petite capitalisation couverte par un
indice à grande capitalisation conserve le spread de taille, qui peut varier de plusieurs
pour cent en un trimestre. Un \omterm{def:m1:futures-contracts-and-exchanges:future}{contrat à terme} sur indice à petite capitalisation couvre mieux et
coûte plus cher à négocier.
\textbf{14.} Une position vendeuse sur \omterm{def:m1:futures-contracts-and-exchanges:future}{contrat à terme} \emph{rapporte} le taux de financement implicite moins
les dividendes~: le fonds reçoit le taux de référence plus le \omterm{def:m1:delta-one-instruments:swap}{spread de financement}
des \omterm{def:m1:futures-contracts-and-exchanges:future}{contrats à terme} sur le notionnel couvert, en échange de la renonciation à la
prime de risque actions.
\textbf{15.} Sélection de titres (voulue), erreur de bêta, décalage de taille et de
secteur, risque de change sur la poche européenne, liquidité pour la \omterm{def:m1:clearing-and-settlement:margin}{marge de variation},
et les heures de chaque jour où les \omterm{def:m1:futures-contracts-and-exchanges:future}{contrats à terme} se négocient et pas les actions.
\textbf{16.} Le \omterm{def:m1:futures-contracts-and-exchanges:future}{contrat à terme} ne peut pas se négocier sous la \omterm{def:m1:asian-and-emerging-equity-markets:limit}{limite de variation de cours} tant que le
marché actions est suspendu~; à leur réouverture, la limite suivante s'applique. La couverture
reste correctement valorisée au règlement, mais ne peut pas être ajustée pendant la suspension, et les
actions ne peuvent pas non plus être vendues.
\textbf{17.} À peine~: le résidu tomberait de 50\,000~\$ à au plus
15\,000~\$, pour une position dont l'incertitude sur le bêta vaut 18.7 millions de dollars. Cela
compte pour des portefeuilles cent fois plus petits.
\textbf{18.} Coût et réversibilité~: vendre 187 millions de dollars d'actions coûte
des dizaines de points de base, réalise des impôts, et abandonne le choix de titres~; 717
\omterm{def:m1:futures-contracts-and-exchanges:future}{contrats à terme} coûtent environ un point de base et laissent le portefeuille intact.
\textbf{19.} \textbf{Vendre 717 E-mini~; résidu $-\$50\,000$.}
\textbf{20.} Il standardise tout sauf le prix, la quantité et
le choix du contrat qui approche le risque que l'on porte réellement.
\end{solution}

\begin{solution}{iq:m1:futures-contracts-and-exchanges:1}
Le multiplicateur vaut 50~\$, le \omterm{def:m1:futures-contracts-and-exchanges:tick}{tick} 0.25 point d'indice, donc un \omterm{def:m1:futures-contracts-and-exchanges:tick}{tick} vaut 12.50~\$. À
un indice de 6\,000, le contrat vaut 300\,000~\$~; un mouvement de 1~\% représente 3\,000~\$ par
contrat.
\iqlookfor{rappel instantané, et réflexe de convertir en notionnel.}
\end{solution}

\begin{solution}{iq:m1:futures-contracts-and-exchanges:2}
Le volume compte les contrats échangés pendant une période~; la position ouverte compte
les contrats existants à un instant donné. Une transaction entre deux nouveaux participants
les ajoute tous deux~; une transaction dans laquelle une partie clôture sa position laisse la position ouverte
inchangée~; une transaction entre deux clôtures la réduit. Un day trader génère
du volume et aucune position ouverte.
\iqlookfor{les trois cas.}
\end{solution}

\begin{solution}{iq:m1:futures-contracts-and-exchanges:3}
Sous forme de \omterm{def:m1:futures-contracts-and-exchanges:tick}{ticks} entiers, avec le pas de cotation et le multiplicateur conservés une seule fois, exactement, dans
les données de référence. Les comparaisons et les clés sont exactes, l'arithmétique est rapide, des prix
hors de la grille ne peuvent pas exister, et le montant monétaire s'obtient par une seule multiplication à
la sortie. Les chaînes décimales comme 0.01 ou les prix en trente-deuxièmes sont
analysées une seule fois à l'entrée. Les pas de cotation changent (et diffèrent pour les spreads et
les options), donc la conversion appartient à l'instrument et à la date.
\iqlookfor{des entiers, un point de conversion unique, et la conscience que les pas de cotation
changent.}
\end{solution}

\begin{solution}{iq:m1:futures-contracts-and-exchanges:4}
Une action est émise par une société et n'importe quelle plateforme peut la négocier~; un \omterm{def:m1:futures-contracts-and-exchanges:future}{contrat à terme} est
créé par une bourse, qui en est propriétaire et le compense dans sa propre
chambre de compensation. La position ouverte n'est pas fongible d'une chambre de compensation à l'autre, si bien qu'un
contrat identique d'un concurrent démarre sans liquidité et sans compensation de marges
avec les autres produits de l'occupant. La liquidité attire alors
la liquidité.
\iqlookfor{la compensation comme verrou, et non la négociation.}
\end{solution}

\begin{solution}{iq:m1:futures-contracts-and-exchanges:5}
Savoir comment la série a été construite~: à quel jour elle change de contrat et
si elle est ajustée. Calculer le rendement de chaque jour à partir de deux cours du
\emph{même} contrat, en changeant de contrat à une date de roll choisie par une règle
(croisement du volume ou de la position ouverte, ou nombre de jours fixe avant l'échéance), de sorte que le
saut de roll n'entre jamais dans un rendement. Choisir entre des niveaux ajustés par ratio et
par différence selon l'usage~; vérifier le règlement-livraison par rapport au dernier
cours échangé~; aligner l'heure du règlement-livraison sur celle de toute autre série utilisée avec elle.
\iqlookfor{des rendements à partir de cours du même contrat~; une règle de roll explicite.}
\end{solution}

\begin{solution}{iq:m1:futures-contracts-and-exchanges:6}
Quel indice~: celui le plus proche du portefeuille, arbitré contre la liquidité.
Bêta~: estimé comment, sur quelle fenêtre, avec quelle erreur. Contrats~: $\beta
V/(mF)$, arrondi. Échéance~: l'\omterm{def:m1:futures-contracts-and-exchanges:code}{échéance la plus proche} et le roll, ou le mois qui couvre
l'horizon. Exécution~: un calendrier au regard de l'urgence de la couverture, peut-être à la
clôture pour coller à la valorisation du portefeuille. Marge~: le dépôt, plus la trésorerie
nécessaire à la \omterm{def:m1:clearing-and-settlement:margin}{marge de variation} lors des hausses. Devise, le cas échéant. Suivi~: dérive du bêta,
expositions résiduelles aux facteurs, le roll. Sortie~: la même chose à l'envers.
\iqlookfor{la trésorerie pour la \omterm{def:m1:clearing-and-settlement:margin}{marge de variation}, que la plupart des candidats oublient.}
\end{solution}
